\chapter{分数：四则运算}\label{ch:g8:fractions}

本章完成\omterm{def:g6:fractions:def}{分数}的全部运算：加法与减法，
\omterm{def:g4:fractions:def}{分母}可以\emph{任意}，乘法，以及——
新来的——除法，其实是一种伪装的乘法。
来自\cref{ch:g8:negprod}的符号法则也一并适用。

\section{分母任意的加法}

\begin{method}[通分]\label{met:g8:fractions:add}
对 $\frac ab$ 与 $\frac cd$ 做加法或减法：
\begin{enumerate}
  \item 求 $b$ 与 $d$ 的一个公\omterm{def:g5:division:multiple}{倍数}（积
        $b \times d$ 总是可行；取较小的公\omterm{def:g5:division:multiple}{倍数}可以省些功夫）；
  \item 把两个\omterm{def:g6:fractions:def}{分数}都化为以该公\omterm{def:g5:division:multiple}{倍数}为\omterm{def:g4:fractions:def}{分母}的\omterm{def:g6:fractions:def}{分数}；
  \item 对\omterm{def:g4:fractions:def}{分子}做加法或减法，并化简。
\end{enumerate}
\end{method}

\begin{example}\label{ex:g8:fractions:add}
计算 $\dfrac{5}{6} + \dfrac{3}{4}$。$6$ 与
$4$ 的一个公\omterm{def:g5:division:multiple}{倍数}是 $12$：
\[
\frac{5}{6} = \frac{10}{12},
\qquad
\frac{3}{4} = \frac{9}{12},
\qquad
\frac{5}{6} + \frac{3}{4} = \frac{10 + 9}{12} = \frac{19}{12}.
\]
含负号时：$\dfrac{2}{3} - \dfrac{7}{5} = \dfrac{10}{15} -
\dfrac{21}{15} = -\dfrac{11}{15}$。
\end{example}

\section{分数的乘法}

\begin{theorem}[分数的乘积]\label{thm:g8:fractions:mult}
\[
\frac{a}{b} \times \frac{c}{d} = \frac{a \times c}{b \times d} .
\]
\end{theorem}

\begin{proof}[从图形看原因]
以 $\frac34$ 乘正方形的 $\frac25$：竖切成 $5$ 条，
保留 $2$ 条；再横切成 $4$ 份，保留剩下部分中的 $3$
行。保留部分是由 $3 \times 2$ 个小格组成的方格网，
它取自 $4 \times 5$ 个小格，即\omterm{def:g6:fractions:def}{分数} $\frac{6}{20}$。
\end{proof}

\begin{omfigure}
\begin{tikzpicture}[scale=0.68]
  \draw[gray!50, thin] (0,0) grid[xstep=1, ystep=1] (5,4);
  \fill[omDef!18] (0,0) rectangle (2,4);
  \fill[omThm!35] (0,0) rectangle (2,3);
  \draw[thick] (0,0) rectangle (5,4);
  \draw[omDef, thick] (2,0) -- (2,4);
  \draw[omThm, thick] (0,3) -- (5,3);
  \node[below, font=\small, omDef] at (1,0) {$\frac25$ 保留};
  \node[right, font=\small, omThm] at (5.05,1.5) {$\frac34$ 的部分};
\end{tikzpicture}

{\small $\frac34 \times \frac25$：双阴影区域是
$3 \times 2 = 6$ 个小格（共 $4 \times 5 = 20$ 个），即\
$\frac{6}{20} = \frac{3}{10}$ ——\omterm{def:g4:fractions:def}{分子}相乘，\omterm{def:g4:fractions:def}{分母}
相乘。}
\end{omfigure}

\begin{example}\label{ex:g8:fractions:mult}
约分要\emph{在相乘之前}进行，划去公\omterm{def:g5:division:multiple}{因数}：
\[
\frac{7}{12} \times \frac{9}{14}
= \frac{7 \times 9}{12 \times 14}
= \frac{7 \times 3 \times 3}{3 \times 4 \times 2 \times 7}
= \frac{3}{8}
\]
（\omterm{def:g4:fractions:def}{分子}、\omterm{def:g4:fractions:def}{分母}中的\omterm{def:g5:division:multiple}{因数} $7$ 与一个 $3$ 被约去）。
含负号时，先由\cref{thm:g8:negprod:rules}确定符号：
$\left(-\frac23\right) \times \left(-\frac{5}{4}\right) =
+\frac{10}{12} = \frac56$。
\end{example}

\section{分数的除法}

\begin{definition}[倒数]\label{def:g8:fractions:inverse}
非零数 $x$ 的\emph{倒数}\index{倒数}是
与 $x$ 相乘得 $1$ 的数。
$\frac ab$ 的倒数是 $\frac ba$，因为
$\frac ab \times \frac ba = \frac{ab}{ab} = 1$；整数 $n$ 的倒数是
$\frac1n$。
\end{definition}

\begin{theorem}[除法即乘以倒数]\label{thm:g8:fractions:div}
设 $c \neq 0$：
\[
\frac{a}{b} \div \frac{c}{d} = \frac{a}{b} \times \frac{d}{c} .
\]
\end{theorem}

\begin{proof}
由定义，商 $q$ 是 $\frac ab$ 除以 $\frac cd$ 所得的数，
即满足 $q \times \frac cd = \frac ab$ 的数。验证候选值
$q = \frac ab \times \frac dc$：
\[
\left(\frac ab \times \frac dc\right) \times \frac cd
= \frac ab \times \left(\frac dc \times \frac cd\right)
= \frac ab \times 1 = \frac ab . \qedhere
\]
\end{proof}

\begin{example}\label{ex:g8:fractions:div}
\[
\frac{3}{5} \div \frac{7}{10}
= \frac{3}{5} \times \frac{10}{7}
= \frac{3 \times 10}{5 \times 7}
= \frac{30}{35} = \frac{6}{7}.
\]
检验一下：除以 $\frac{7}{10}$（一个小于
$1$ 的数）的结果必须\emph{大于} $\frac35$ ——而且
事实上 $\frac67 > \frac35$（$\frac{30}{35} > \frac{21}{35}$）。
\end{example}

\begin{example}[分数线中的分数线]\label{ex:g8:fractions:complex}
“双层\omterm{def:g6:fractions:def}{分数}”其实就是竖着写的除法：
\[
\frac{\ \frac{2}{3}\ }{\ \frac{5}{6}\ }
= \frac{2}{3} \div \frac{5}{6}
= \frac{2}{3} \times \frac{6}{5}
= \frac{12}{15} = \frac{4}{5}.
\]
先找出\emph{主}\omterm{def:g6:fractions:def}{分数}线（即最长的那条），再
应用\cref{thm:g8:fractions:div}。
\end{example}

\begin{method}[混合运算]\label{met:g8:fractions:mixed}
在含有\omterm{def:g6:fractions:def}{分数}与四则运算的式子中：
\begin{enumerate}
  \item 按通常的运算顺序进行（见\cref{ch:g7:priorities}）；
  \item 将每个除法都化为乘以\omterm{def:g8:fractions:inverse}{倒数}；
  \item 先确定符号（见\cref{met:g8:negprod:compute}），再
        相乘或通分；
  \item 尽早约分。
\end{enumerate}
\end{method}

\begin{example}\label{ex:g8:fractions:mixed}
\begin{align*}
\frac12 + \frac{2}{3} \times \frac{9}{4}
&= \frac12 + \frac{2 \times 9}{3 \times 4}
&& \text{（先算乘积！）} \\
&= \frac12 + \frac{18}{12} = \frac12 + \frac32
= \frac{4}{2} = 2 .
\end{align*}
\end{example}

\section{练习}

\begin{exercise}[$\star$]\label{exo:g8:fractions:1}
计算并化简：
\[
\frac{3}{4} + \frac{2}{5}, \qquad
\frac{5}{6} - \frac{3}{8}, \qquad
\frac{7}{10} + \frac{8}{15} .
\]
\end{exercise}

\begin{exercise}[$\star$]\label{exo:g8:fractions:2}
计算（先确定符号）：
\[
-\frac{2}{7} + \frac{3}{14}, \qquad
\frac{1}{6} - \frac{5}{4}, \qquad
-\frac{3}{5} - \frac{1}{2} .
\]
\end{exercise}

\begin{exercise}[$\star$]\label{exo:g8:fractions:3}
计算，相乘前先约分：
\[
\frac{5}{8} \times \frac{4}{15}, \qquad
\frac{9}{14} \times \frac{7}{3}, \qquad
\left(-\frac{6}{5}\right) \times \frac{10}{9} .
\]
\end{exercise}

\begin{exercise}[$\star$]\label{exo:g8:fractions:4}
写出下列各数的\omterm{def:g8:fractions:inverse}{倒数}：$\dfrac{3}{7}$；\quad $5$；\quad
$\dfrac{1}{4}$；\quad $-\dfrac{2}{9}$。
\end{exercise}

\begin{exercise}[$\star$]\label{exo:g8:fractions:5}
计算：
\[
\frac{4}{9} \div \frac{2}{3}, \qquad
\frac{7}{5} \div 14, \qquad
6 \div \frac{3}{4} .
\]
\end{exercise}

\begin{exercise}[$\star$]\label{exo:g8:fractions:6}
计算下列“双层\omterm{def:g6:fractions:def}{分数}”：
\[
\frac{\ \frac{3}{4}\ }{\ \frac{9}{8}\ },
\qquad
\frac{\ \frac{5}{6}\ }{\ 10\ },
\qquad
\frac{\ 2\ }{\ \frac{4}{7}\ } .
\]
\end{exercise}

\begin{exercise}[$\star$]\label{exo:g8:fractions:7}
计算，注意运算顺序：
\[
\frac{1}{3} + \frac{1}{2} \times \frac{4}{5},
\qquad
\left(\frac{1}{3} + \frac{1}{2}\right) \times \frac{4}{5} .
\]
\end{exercise}

\begin{exercise}[$\star\star$]\label{exo:g8:fractions:8}
一个水箱装了 $\frac{2}{5}$，再加入\emph{水箱\omterm{def:g2:measure:units}{容积}}的 $\frac{1}{3}$。
此时水箱中有几分之几是满的？
还有几分之几是空的？
\end{exercise}

\begin{exercise}[$\star\star$]\label{exo:g8:fractions:9}
某班有四分之三的学生通过了测验；在这些学生中，
三分之二的人得分超过 $14$（满分 $20$）。全班中得分超过 $14$ 的
人占几分之几？若这部分人数为 $12$ 人，则
全班共有多少人？
\end{exercise}

\begin{exercise}[$\star\star$]\label{exo:g8:fractions:10}
一根长 $\frac{15}{2}$ 米的绳子，要剪成每段 $\frac{3}{4}$ 米的
小段。可以剪得几段？（这是\omterm{def:g6:fractions:def}{分数}除法——
请检验结果是否为整数。）
\end{exercise}

\begin{exercise}[$\star\star$]\label{exo:g8:fractions:11}
计算
\[
E = \frac{2}{3} - \frac{5}{6} \div \frac{10}{9},
\qquad
F = \left(-\frac{1}{2}\right)^2 \times \frac{8}{3} .
\]
\end{exercise}

\begin{exercise}[$\star\star\star$]\label{exo:g8:fractions:12}
化简下列表达式
\[
1 + \cfrac{1}{1 + \cfrac{1}{1 + 1}}
\]
（从最内层的\omterm{def:g6:fractions:def}{分数}开始，由内向外，
每次处理一条\omterm{def:g6:fractions:def}{分数}线）。
\end{exercise}

\section{问题：埃及分数}

\begin{problem}[{周末问题——每个分数都可写成
互不相同的单位分数之和}]
\label{pb:g8:fractions:1}
\omterm{def:g4:fractions:def}{分子}为 $1$ 的\omterm{def:g6:fractions:def}{分数}称为\emph{单位\omterm{def:g6:fractions:def}{分数}}。
古埃及的抄写员——莱因德纸草书大约抄写于
公元前 $1550$ 年——把每个\omterm{def:g6:fractions:def}{分数}都写成\emph{互不相同}的
单位\omterm{def:g6:fractions:def}{分数}之和：绝不写成 $\frac27 = \frac17 + \frac17$，而是
写成 $\frac27 = \frac14 + \frac{1}{28}$。这样的和式称为该\omterm{def:g6:fractions:def}{分数}的
\emph{埃及记法}。本题将
本章的加法与比较技巧
（\cref{met:g8:fractions:add}）发展为一种方法——该方法由斐波那契于
$1202$ 年发表，能将任一
$0$ 与 $1$ 之间的\omterm{def:g6:fractions:def}{分数}写成埃及记法，并以一道著名的
遗产分配难题作结。

\medskip
\noindent\textbf{第一部分——单位\omterm{def:g6:fractions:def}{分数}与拆分
恒等式。}
\begin{enumerate}
  \item 计算 $\frac12 + \frac13$ 与 $\frac12 + \frac14$，并
        由此写出 $\frac56$ 与 $\frac34$ 的埃及记法。
  \item 验证抄写员对 $\frac27$ 的记载：说明
        $\frac14 + \frac{1}{28} = \frac27$。
  \item 检验 $\frac13 + \frac16 = \frac12$，以及
        $\frac14 + \frac{1}{12} = \frac13$。然后证明
        \emph{拆分恒等式}：对每个正整数
        $n \geq 1$，有
        \[
        \frac{1}{n} = \frac{1}{n+1} + \frac{1}{n(n+1)} .
        \]
  \item 由此写出 $1$ 本身的埃及记法，表示为
        三个互不相同的单位\omterm{def:g6:fractions:def}{分数}之和。
  \item 将等式 $\frac56 = \frac12 + \frac13$ 中的最后一项拆开，
        得到 $\frac56$ 的\emph{第二种}埃及记法，含
        三个单位\omterm{def:g6:fractions:def}{分数}。说明为什么任何\omterm{def:g6:fractions:def}{分数}的
        埃及记法都不是唯一的。
\end{enumerate}

\medskip
\noindent\textbf{第二部分——斐波那契的贪心法。}
该方法为：每次减去可以减去的\emph{最大}单位\omterm{def:g6:fractions:def}{分数}，
并对剩余部分重复这一过程。
\begin{enumerate}[resume]
  \item 利用通分，说明
        $\frac14 < \frac27 < \frac13$。为什么这能证明
        $\frac14$ 是小于
        $\frac27$ 的最大单位\omterm{def:g6:fractions:def}{分数}？
  \item 计算 $\frac27 - \frac14$，并指出
        贪心法给出的是 $\frac27$ 的哪一种记法？
  \item 取 $\frac57$，说明
        小于 $\frac57$ 的最大单位\omterm{def:g6:fractions:def}{分数}是 $\frac12$，并计算
        $\frac57 - \frac12$。
  \item 说明小于
        $\frac{3}{14}$ 的最大单位\omterm{def:g6:fractions:def}{分数}是 $\frac15$（将 $\frac{3}{14}$ 与
        $\frac14$、$\frac15$ 比较），计算
        $\frac{3}{14} - \frac15$，并得出：
        \[
        \frac57 = \frac12 + \frac15 + \frac{1}{70} .
        \]
        用公\omterm{def:g4:fractions:def}{分母}直接验证这一记法，
        公\omterm{def:g4:fractions:def}{分母}为 $70$。
  \item 对 $\frac45$ 运行贪心法，每一步检验
        所取的单位\omterm{def:g6:fractions:def}{分数}是否为可以减去的最大者，并
        验证最终的记法。
\end{enumerate}

\medskip
\noindent\textbf{第三部分——为何该方法总会停止，以及
第十八头骆驼。}
\begin{enumerate}[resume]
  \item 在第 7 至 10 题中，该方法分别作用于 $\frac27$、
        $\frac57$ 与 $\frac45$。对每个\omterm{def:g6:fractions:def}{分数}，列出
        各次\omterm{def:g3:division:remainder}{余数}的\omterm{def:g4:fractions:def}{分子}（化简之前）。
        你观察到什么？
  \item 设\omterm{def:g6:fractions:def}{分数} $\frac ab$ 与正整数 $n$，
        证明下式成立：
        \[
        \frac{a}{b} - \frac{1}{n}
        = \frac{a \times n - b}{b \times n} .
        \]
  \item 设贪心法减去的是 $\frac1n$，
        即从 $\frac ab$ 中减去它：于是 $\frac1n$ 可以减去，但 $\frac{1}{n-1}$
        已经太大，$\frac{1}{n-1} > \frac ab$。将两者
        以 $b \times (n - 1)$ 为公\omterm{def:g4:fractions:def}{分母}通分，推出
        $b > a \times n - a$，从而新\omterm{def:g4:fractions:def}{分子}
        满足
        \[
        a \times n - b < a .
        \]
        说明为何这证明了该方法总会停止。
  \item 一个古老的故事：一位父亲去世，留下 $17$ 头骆驼，遗嘱
        将骆驼群的 $\frac12$ 给长子，$\frac13$ 给
        次子，$\frac19$ 给三子——但每份都
        不是整数头骆驼。一位邻居借给这个家庭
        第十八头骆驼；孩子们分别取走 $9$、$6$ 和
        $2$ 头，借来的骆驼随后归还。计算
        $\frac12 + \frac13 + \frac19$，并解释其中的诀窍：为何
        骆驼群为何能按整数头分配，以及为何
        没有孩子所得少于遗嘱承诺的份额？
  \item 修正遗嘱：提出三个\emph{互不相同}的单位\omterm{def:g6:fractions:def}{分数}，
        其和恰好为 $1$，并给出
        使三份均为整数头骆驼的最小
        骆驼群规模。
\end{enumerate}
\end{problem}
